\documentclass[tikz,border=2bp]{standalone}
\usepackage{ctex}
\usepackage[europeanresistors, americanvoltages, americancurrents, siunitx]{circuitikz}
\standaloneenv{circuitikz}

\begin{document}
\begin{circuitikz}[semithick, scale=1.0, transform shape]

  % 1A电流源
  \draw (0,0) to[isource, l={\SI{1}{\ampere}}] (0,3)
        -- (2.0,3);
        
  % 第一支路：6Ω电阻 (使用node手动精准标在左侧)
  \draw (2.0,3) to[R, *-*] (2.0,0);
  \node[left] at (1.5,1.5) {$6\text{ }\Omega$};
  
  % 第二支路：开关S(t=0闭合) 与 6Ω电阻串联 (手动控制左侧文字定位，右侧电流定位)
  \draw (2.0,3) -- (4.5,3)
        to[cspst] (4.5,1.8)
        to[R, i_>={$i(t)$}] (4.5,0);
  \node[left] at (4.0,2.4) {$t=0$};
  \node[left] at (4.0,0.9) {$6\text{ }\Omega$};
        
  % 第三支路：电容C=1F (手动控制左侧容量定位，右侧电压定位)
  \draw (4.5,3) -- (7.0,3)
        to[C] (7.0,0);
  \node[left] at (6.3,1.5) {$C = \SI{1}{\farad}$}; % x移至 6.3 彻底消除与电容极板的碰撞
  \node[right] at (7.2,1.9) {$+$};
  \node[right] at (7.2,1.1) {$-$};
  \node[right] at (7.2,1.5) {$u_C(t)$};
        
  % 底部公共导线与接地
  \draw (7.0,0) -- (0,0);
  \draw (3.25,0) node[ground] {};

\end{circuitikz}
\end{document}
